\section{问题与成功标准}
\label{sec:problem}

\paragraph{分析问题}
胃癌原发灶的空间切片上，成纤维高分区和巨噬高分区是否比随机标签更常落在150\,$\mu$m以内；若以界面上的配体受体乘积作为虚拟阻断的读数，哪些边在坐标打乱之后仍然更高。腹膜单细胞只用来看这些基因在腹膜肿瘤和原发肿瘤之间的样本均值方向。

\paragraph{成功标准}
\begin{itemize}
  \item 主指标：界面信号。对一条配体受体，在巨噬高分区且150\,$\mu$m内至少有一个成纤维高分区邻居的点位上，计算受体表达与这些邻居配体均值的乘积，再取中位数。表达为每1万计数的$\log(1+x)$。
  \item 接触指标：巨噬高分区点位中，150\,$\mu$m内存在成纤维高分区邻居的比例。
  \item 对照：接触指标打乱两类标签并保持数目，200次，取第95百分位。界面信号打乱点位坐标，200次，取第95百分位。种子为20261002。
  \item 保留规则：原发灶9例中，至少5例的界面信号高于该切片空分布，并且去掉任意一例后，其余8例里仍至少有5例高于空分布。只达到5例、留一后掉到4例的边不保留。
\end{itemize}

接触比例的第95百分位和界面信号的5/9例规则，是本分析在出结果之前写下的项目阈值，不是外部标准。GC6-PM不进入上述计数。

\paragraph{范围}
覆盖GSE251950的10张Visium切片和GSE183904中标成腹膜肿瘤或原发肿瘤的样本。不覆盖腹膜灶上的空间虚拟阻断，也不把样本均值方向写成治疗排序。

\paragraph{事后指定的腹膜乘积}
类均值看完之后，另做一次样本水平的乘积比较。配体和受体各自取细胞类均值，再相乘。只检验下面7条的腹膜中位数是否高于原发：C3成纤维或上皮乘C3AR1巨噬，SPP1巨噬乘CD44巨噬或T细胞，CCL2上皮乘CCR2 T细胞，CCL2成纤维乘CCR2巨噬，TGFB1 T细胞乘TGFBR2巨噬。对照是在有限值的样本里穷举哪3份标成腹膜。这不是空间分析之前写死的主终点，也不把细胞个数当作样本量。
